\chapter{Dry Storage and Longevity of Seeds}

Ultimately dry storage is fatal to all seeds including those of the D-70 type even
though these require dry storage in order to germinate. The dieing of seeds in dry
storage is a rate process so that no single time can be given for the retention of
viability. It will also depend on temperature and relative humidity. There is also the
problem that dry stored seeds may still germinate, but the seedlings may be weak and
unable to survive because of internal deterioration in the seeds.

The longevity of seeds in dry storage has been the subject of much interest
(refs. 2-7), and an entire book has been devoted to this subject (ref. 6). There are
studies on the viability of seeds in dry storage that have been conducted continuously
since the 1850-1860 period (refs. 6 and 33). The results can be summarized
succinctly. A few species have seed that retain viability for only days such as Populus,
Salix, and Tussilago. Most species retain viability for a year or two. Some species
retain viability for 10-50 years. Some species of Fabaceae have seeds that are still
viable after 100-150 years in dry storage and a few Malvaceae are also long lived (ref.
33). The seeds with long lives are invariably seeds with impervious seed coats.

Tabloid stories claiming viability of ancient seeds are false, even though they
have been repeated in reputable journals. Examples of such stories are reports of
viable seeds of grains being found clutched in the hands of Egyptian mummies, viable
seeds of lotus that are ten thousand years old being dug up in ancient peat beds in
China, and viable seeds that are thousands of years old being found in the caches of
Arctic rodents. All of these have been disproven (refs. 2-6).

In the first few years of the program it was the standard procedure to compare
the behavior of fresh seeds with seeds dry stored for six months at 40 and seeds dry
stored for six months at 70. When it was found that differences were small: the
experiments on dry storage at 40 were discontinued. The data in Table 13-1 illustrate
these small differences. The practice of storing seeds in the refrigerator is overrated.

\begin{table}[h!]
\refstepcounter{table}

\begin{center}
Table \thetable. Percent Germinations for Seeds Dry Stored (DS) for Six Months
at 40 Relative To 70
\end{center}

\begin{tabular}{|l|l|p{2cm}|l|}
\hline
Species & Fresh Seed & DS at 40 & DS at 70 \\
\hline
Hyacinthus orientalis & 70-40(100\%)      & 70(85\% rotted)-40-70-40-70-40(11\%) & 70(100\% rotted) \\
Scilla sibirica       & 70(22\%)-40(65\%) & 70-40(53\%)-70(12\%)                 & 70(100\% rotted) \\
Trillium albidum      & 70(74\%)-40(11\%) & 70-40(20\%)                          & 70(100\% rotted) \\
\hline
\end{tabular}
\end{table}

Seeds of many species are totally killed by dry storage for six months at either
40 or 70. Some examples are Anemonella thalictroides, Chelidonium majus,
Claytonia virginica, Corydalis cheilanthifolia, Delphinium tricorne, Dentaria laciniata,
Dicentra cucuharia and spectabilis, Jeffersonis diphylla and dubia, Leucojum vernum,
Meconopsis cambrica, Populus tremuloides, Ranunculus hispidulus, Salix,
Sanguinaria canadensis, Stylophorum diphyllum, Tiarella cordifolia, Trillium
grandiflorum, and Tussilago farfara.

Dry storage is not the only way to store seeds. The germination pattern of many
species offers the opportunity for storing the seeds in a moist condition, and this is the
answer for seeds that cannot tolerate dry storage. For example, a 40-70 germinator
such as Nemastylis acuta could be stored moist at 70 for many months and possibly a
year or more. This period of time in moisture at 70 is a true dormant period, and
destruction of the germination inhibitors will not not begin until the shift to 40.
With a 70-40 germinator like Eranthis hyemalis, the seeds could be stored moist
at 40 for many months and possibly a year or more. This period of time in moisture at
40 is a true dormant period and destruction of the germination inhibitors will not begin
until a shift to 70. Alternatively seeds of Eranthis could be stored moist at 70 for at
least six months because even though inhibitors are being destroyed in the first three
months under these conditions, germination will not take place until the shift to 40.

Storing seeds in moist towels in polyethylene bags is nearly as convenient as
storing seeds in paper envelopes, and seeds stored in this manner can be shipped
anywhere. I can confidently predict that such a practice will become widespread. The
information presented in Chapter 20 is sufficient to select the species which have
germination patterns adapted to such procedures.

A major question is whether the inhibitor destruction reactions require oxygen.
Using Eranthis hyemalis as an example, if the inhibitor destruction reactions did not
require oxygen, the seeds could be stored and distributed at 70 in sealed polyethylene
packets containing a drop of water. If the reactions required only small amounts of
oxygen, possibly the diffusion of oxygen through the polyethylene would be sufficient,
and again the seeds could be placed in sealed packets. However, if the inhibitor
destruction reactions consumed large amounts of oxygen, some provision would have
to be made for keeping the packet of moist seeds aerated. Perhaps a small hole in the
polyethylene would be sufficient. All of these questions could be readily answered
with a few direct experiments, but the answers might well vary from species to species.

Another possibility that deserves more study is the storage of seeds moist in the
presence of germination inhibitors. Seeds in fruits are a natural example. The
procedures could be refined by adding preservatives to stop fermentation and by
refrigeration. Abscicic acid blocks the germination in many species (ref. 24), and
perhaps this could be used to delay germination in seeds being stored moist.
There have been frequent reports that seeds stored in the presence of drying
agents such as silica gel retain viability better than seed stored under conditions of
higher humidity (ref. 34). This would only apply to seeds that survive dry storage. To
do this sort of work correctly, the phenomena have to be treated as rates and the rate
characteristics measured. Reports that seeds store better in sealed containers
appears to be only due to lower humidity (ref. 35). Experiments were conducted on
dry storage of seeds at 40 in and out of polyethylene bags. These were discontinued
when no differences were found.

There have been studies on the preservation of seeds in liquid nitrogen (ref.
36). Special precautions were necessary to get the moisture content of the seed
unusually low. At best this type of storage would be suitable only for seed that can be
dry stored, and inferences that it has general applicability are false.

